\subsection{Generation and logical translations}
\label{sec:internal-adjunctions}

Continue with $\mathscr S=\Split(\mathscr R)$ and
$\mathcal D=\Tot(\mathscr S)$. Assume that $\mathcal D$ is regular
with disjoint, pullback-stable finite coproducts. Fix a first-order
syntax as in Definition~\ref{def:categorical-syntax}. For every
axiom subobject of its sentence object, assume an initial algebra
of derivations as in Appendix~\ref{app:first-order-derivations}. For an object $X$ of $\mathcal D$, write
$\operatorname{Sub}_{\mathcal D}(X)$ for the ordered collection of
subobjects of $X$, with order given by factorization. If
$f:X\to Y$ is total, write $f^*B$ for the inverse image of
$B\hookrightarrow Y$, and $f_!A$ for the image of
$A\hookrightarrow X\xrightarrow{f}Y$.

Let $j:S\hookrightarrow\Phi$ be the sentence inclusion. For
$A\hookrightarrow S$, let $C(A)\hookrightarrow\Phi$ be the least
closed formula subobject supplied by
Theorem~\ref{thm:first-order-theory}, and put
\[
 \operatorname{Gen}(A)=j^*C(A)=T_{A,S}.
\]
Write $\operatorname{Th}_{\mathcal D}(S)$ for the sentence theories,
ordered by inclusion. Equivalently, these are the subobjects $T$ with
$\operatorname{Gen}(T)=T$.

\begin{theorem}[Generation reflection]\label{thm:first-order-reflection}
Under the preceding hypotheses, $\operatorname{Gen}$ is extensive,
monotone, and idempotent. It is left adjoint to the inclusion
$\iota:\operatorname{Th}_{\mathcal D}(S)\hookrightarrow
\operatorname{Sub}_{\mathcal D}(S)$. Thus, for every axiom
subobject $A$ and sentence theory $T$,
\[
 \operatorname{Gen}(A)\leq T\quad\Longleftrightarrow\quad A\leq T.
\]
For each theory $T$ choose a representative $i_T:|T|\hookrightarrow S$.
Let $J:\operatorname{Th}_{\mathcal D}(S)\to\mathcal D/S$ send $T$
to $i_T$ and $T\leq U$ to the unique factorization of $i_T$ through $i_U$. Then
$\Gamma(a:A\to S)=\operatorname{Gen}(\im a)$ defines a functor
$\Gamma:\mathcal D/S\to\operatorname{Th}_{\mathcal D}(S)$ with
$\Gamma\dashv J$.
\end{theorem}
\begin{proof}
Axiom introduction gives $A\leq\operatorname{Gen}(A)$, and
leastness of $C(A)$ gives monotonicity. The closed formula
subobject $C(A)$ contains $j_!\operatorname{Gen}(A)$, whence
$C(\operatorname{Gen}(A))\leq C(A)$. Pullback along $j$ and
extensivity prove idempotence. For a theory $T$, monotonicity and $\operatorname{Gen}(T)=T$ give
$\operatorname{Gen}(A)\leq T$ if $A\leq T$. Extensivity gives the converse. If $k:(a:A\to S)\to(b:B\to S)$ lies over $S$, then $a=bk$,
so $\im a\leq\im b$. Applying generation gives $\Gamma(k)$.
A morphism in $\mathcal D/S$ from $a$ to $J(T)$ exists uniquely
exactly when $\im a\leq T$.
The equivalence identifies this hom-set with the hom-set from
$\Gamma(a)$ to $T$ in the theory order, naturally in both objects.
\end{proof}

A first-order presentation $P$ consists of the named-variable
syntax of Definition~\ref{def:categorical-syntax}, the logical rules of
Appendix~\ref{app:first-order-derivations}, and a choice of initial
derivation algebra for each subobject $A\hookrightarrow S_P$ of its
sentence object. Choosing $A$ specifies the nonlogical axioms of a
theory in $P$. Write $\Tm_P,\Phi_P,S_P,j_P,\nu_P$ for the term
object, formula object, sentence object, sentence inclusion, and
sentence negation. Write $\operatorname{Gen}_P$ for generation.

Fix a common variable object $V$ for two such presentations $P,Q$.
A logical translation $\sigma:P\to Q$ consists of total morphisms
\[
 \sigma_{\Tm}:\Tm_P\to\Tm_Q,\qquad
 \sigma_{\Phi}:\Phi_P\to\Phi_Q,\qquad
 \sigma_S:S_P\to S_Q.
\]
The term morphism fixes variables, $\sigma_{\Tm}\eta_P=\eta_Q$,
where $\eta_P:V\to\Tm_P$ and $\eta_Q:V\to\Tm_Q$ are the variable constructors.
The formula morphism preserves equality, negation, implication, and quantification by the equations
\begin{align*}
 \sigma_\Phi\mathsf{eq}_P&=\mathsf{eq}_Q(\sigma_{\Tm}\times\sigma_{\Tm}),&
 \sigma_\Phi\neg_P&=\neg_Q\sigma_\Phi,\\
 \sigma_\Phi({\Rightarrow}_P)&=({\Rightarrow}_Q)(\sigma_\Phi\times\sigma_\Phi),&
 \sigma_\Phi{\forall}_P&={\forall}_Q(1_V\times\sigma_\Phi).
\end{align*}
The sentence morphism satisfies $j_Q\sigma_S=\sigma_\Phi j_P$.
Since $j_Q$ is monic and formula negation restricts to sentence
negation in each presentation, $\sigma_S\nu_P=\nu_Q\sigma_S$.

For substitution, put $E_P=V\times\Tm_P$, $E_Q=V\times\Tm_Q$ and
$e=1_V\times\sigma_{\Tm}:E_P\to E_Q$. Write
$a_{T,P}:\Tm_P\times E_P\to\Tm_P$ for term substitution and
$i_P:\Adm_P\hookrightarrow\Phi_P\times E_P$,
$a_{F,P}:\Adm_P\to\Phi_P$ for the admissibility inclusion and
formula substitution. Use the corresponding notation for $Q$.
We require $\sigma_{\Tm}a_{T,P}=a_{T,Q}(\sigma_{\Tm}\times e)$
and a morphism $s:\Adm_P\to\Adm_Q$ satisfying
\[
 i_Q s=(\sigma_\Phi\times e)i_P,\qquad
 a_{F,Q}s=\sigma_\Phi a_{F,P}.
\]
Thus each admissible substitution in $P$ translates to an admissible
substitution in $Q$ with the indicated value.

For $n\in\{0,1,2\}$, write $R_n^P$ for the object of logical
rule instances with $n$ premises, with premise and conclusion morphisms
\[
 p_n^P:R_n^P\to\Phi_P^n,\qquad q_n^P:R_n^P\to\Phi_P,
\]
and use the corresponding notation for $Q$.
Preservation of the logical rules means that there are morphisms
$\rho_n:R_n^P\to R_n^Q$ with
\[
 p_n^Q\rho_n=\sigma_{\Phi}^{n}p_n^P,\qquad
 q_n^Q\rho_n=\sigma_{\Phi}q_n^P.
\]
For $n=0$, $\rho_0$ sends instances of the seven logical axiom
schemes to logical axiom instances. For generalization and modus ponens
one may take $\rho_1=1_V\times\sigma_\Phi$ and
$\rho_2=\sigma_\Phi\times\sigma_\Phi$, respectively.
The three morphisms $\sigma_{\Tm},\sigma_\Phi,\sigma_S$ determine the
translation. The conditions on $s$ and $\rho_n$ require their existence.
Composition is componentwise. A nonlogical predicate
may be interpreted by a formula. Put
$\operatorname{Th}_P=\operatorname{Th}_{\mathcal D}(S_P)$.

\begin{theorem}[Translation Galois connection]\label{thm:translation-galois}
Let $\sigma:P\to Q$ be a logical translation between presentations
satisfying the hypotheses of
Theorem~\ref{thm:first-order-reflection}. Inverse image along
$\sigma_S$ sends sentence theories to sentence theories, and
\[
 \sigma_\#:\operatorname{Th}_P\rightleftarrows
 \operatorname{Th}_Q:\sigma^*,\qquad
 \sigma_\#(T)=\operatorname{Gen}_Q(\sigma_!T),\qquad
 \sigma_\#\dashv\sigma^*.
\]
Here $\sigma_!$ and $\sigma^*$ abbreviate the image and inverse-image
operations for $\sigma_S$. Moreover,
\[
 \operatorname{Gen}_Q\sigma_!
   =\sigma_\#\operatorname{Gen}_P.
\]
For composable logical translations $\sigma:P\to Q$ and
$\rho:Q\to R$,
\[
 (\rho\sigma)_\#=\rho_\#\sigma_\#,\qquad
 (\rho\sigma)^*=\sigma^*\rho^*.
\]
\end{theorem}
\begin{proof}
Let $T$ be a $Q$-theory. Put $D=C_Q(T)\hookrightarrow\Phi_Q$ and
$E=\sigma_\Phi^*D\hookrightarrow\Phi_P$. Closure of $D$ under the
$n$-premise rule means $(p_n^Q)^*(D^n)\leq(q_n^Q)^*D$ in
$\operatorname{Sub}_{\mathcal D}(R_n^Q)$. Pulling back along $\rho_n$
gives $(p_n^P)^*(E^n)\leq(q_n^P)^*E$. For $n=0$, the left side is
the whole object of logical axiom instances. Since
$j_P^*E=\sigma_S^* j_Q^*D=\sigma^*T$, also $j_{P,!}\sigma^*T\leq E$. Thus
$C_P(\sigma^*T)\leq\sigma_{\Phi}^*C_Q(T)$.
Pullback along $j_P$ gives
$\operatorname{Gen}_P(\sigma^*T)\leq\sigma^*T$;
extensivity gives equality. Now image and inverse image satisfy
$\sigma_!A\leq B$ exactly when $A\leq\sigma^*B$.
Combining this with Theorem~\ref{thm:first-order-reflection}
proves the adjunction.

Apply closure preservation to
$T=\operatorname{Gen}_Q(\sigma_!A)$. The resulting inclusion
$\sigma_!\operatorname{Gen}_P(A)
\leq\operatorname{Gen}_Q(\sigma_!A)$, together with extensivity
and idempotence, proves the generation identity. Composition of
images and pullbacks then gives the two composition identities.
\end{proof}

Let $\mathsf{Ax}_V(\mathcal D)$ have pairs $(P,A)$ as objects,
where $P$ has variable object $V$ and $A\hookrightarrow S_P$ is an
axiom subobject. Its morphisms are logical translations
$\sigma:(P,A)\to(Q,B)$ satisfying $\sigma_!A\leq B$.
Let $\mathsf{Th}_V(\mathcal D)$ be its full subcategory on sentence
theories. Theorem~\ref{thm:translation-galois} also gives a
reflection
\[
 \mathsf{Ax}_V(\mathcal D)\rightleftarrows
 \mathsf{Th}_V(\mathcal D),\qquad
 (P,A)\longmapsto(P,\operatorname{Gen}_P(A)).
\]
The right adjoint is inclusion. The hom-set bijection leaves the
translation unchanged, since, for a target theory $T$,
$\sigma_!A\leq T$ exactly when
$\sigma_!\operatorname{Gen}_P(A)\leq T$.

\subsection{Change of the ambient category}
\label{sec:deduction-change-of-base}

Let $\mathscr S_i$ be split restriction categories and put
$\mathcal D_i=\Tot(\mathscr S_i)$ for $i=1,2$. Assume the
generation hypotheses above for a presentation in each
$\mathcal D_i$, with sentence object $S_i$. Write
$\operatorname{Th}_i=\operatorname{Th}_{\mathcal D_i}(S_i)$ and
let $\Gamma_i\dashv J_i$ be the functors of
Theorem~\ref{thm:first-order-reflection}. Let
$F:\mathscr S_1\to\mathscr S_2$ preserve restrictions and induce
a regular functor $G:\mathcal D_1\to\mathcal D_2$. A regular
functor preserves finite limits and regular epimorphisms. Fix a
total morphism $\sigma:GS_1\to S_2$ and define
\[
 L(A)=\im(\sigma G i_A),\qquad
 K(a:A\to S_1)=\sigma Ga:GA\to S_2.
\]
Here $i_A:A\hookrightarrow S_1$ represents the subobject $A$.
Thus $L:\operatorname{Sub}_{\mathcal D_1}(S_1)\to
\operatorname{Sub}_{\mathcal D_2}(S_2)$ is monotone. The functor
$K:\mathcal D_1/S_1\to\mathcal D_2/S_2$ sends a morphism $k$ over $S_1$
to $Gk$ over $S_2$.
For each $A$, let $c_{A,S_1}:P_{A,S_1}\to S_1$ be the pullback
of its conclusion morphism $c_A:P_A\to\Phi_1$ along the sentence
inclusion $j_1:S_1\hookrightarrow\Phi_1$.

\begin{proposition}[Comparison of deduction]\label{prop:deduction-comparison}
With these data, suppose that for every $A\hookrightarrow S_1$
there is a total morphism
\[
 p_A:GP_{A,S_1}\to P_{L(A),S_2},\qquad
 c_{L(A),S_2}p_A=\sigma Gc_{A,S_1}.
\]
Then $L^{\mathrm{th}}(T)=\operatorname{Gen}_2L(T)$ defines a
functor on theory orders satisfying
\[
 L^{\mathrm{th}}\operatorname{Gen}_1=\operatorname{Gen}_2L,
 \qquad \Gamma_2K=L^{\mathrm{th}}\Gamma_1.
\]
Consequently the left-adjoint square in
\[
\begin{tikzcd}[column sep=large,row sep=large]
 \mathcal D_1/S_1\arrow[r,"K"]
   \arrow[d,shift right=.8ex,"\Gamma_1"'] &
 \mathcal D_2/S_2\arrow[d,shift right=.8ex,"\Gamma_2"']\\
 \operatorname{Th}_1\arrow[r,"L^{\mathrm{th}}"']
   \arrow[u,shift right=.8ex,"J_1"'] &
 \operatorname{Th}_2\arrow[u,shift right=.8ex,"J_2"']
\end{tikzcd}
\qquad \Gamma_i\dashv J_i
\]
commutes. Factorization through the generated theory gives a natural
transformation $KJ_1\to J_2L^{\mathrm{th}}$.
If $L\dashv R$ on subobject orders, $R$ sends theories to theories
and $L^{\mathrm{th}}\dashv R|_{\operatorname{Th}_2}$.
\end{proposition}
\begin{proof}
Factor $c_{A,S_1}$ through its image $\operatorname{Gen}_1(A)$.
Since $G$ preserves the regular epimorphism in this factorization,
\[
 L(\operatorname{Gen}_1(A))
 =\im(\sigma Gc_{A,S_1})
 =\im(c_{L(A),S_2}p_A)
 \leq\operatorname{Gen}_2(L(A)).
\]
Applying $\operatorname{Gen}_2$ gives one inequality in the first
identity by idempotence. The other follows from
$A\leq\operatorname{Gen}_1(A)$ by monotonicity of $L$ and
$\operatorname{Gen}_2$. Factoring $a$ through its image in the same
way gives $\im(K(a))=L(\im a)$, proving the second identity.
For each theory $T$, the unique factorization of $KJ_1(T)$ through
$J_2L^{\mathrm{th}}(T)$ defines the component of the comparison at $T$. Uniqueness proves naturality.
For a closed $B$, adjunction gives $LRB\leq B$, and hence
\[
 L\operatorname{Gen}_1(RB)
 \leq\operatorname{Gen}_2(LRB)\leq B.
\]
Thus $\operatorname{Gen}_1(RB)\leq RB$. Extensivity gives
equality, and the restricted adjunction follows from the
generation reflection.
\end{proof}

Suppose $G$ has a right adjoint $G^{\mathrm R}:\mathcal D_2\to\mathcal D_1$
with unit $\eta:1_{\mathcal D_1}\to G^{\mathrm R}G$. Put
$\sigma^\flat=G^{\mathrm R}\sigma\,\eta_{S_1}:S_1\to G^{\mathrm R}S_2$.
For $i_B:B\hookrightarrow S_2$, let $R(B)=(\sigma^\flat)^*(G^{\mathrm R} i_B)$.
The right adjoint $G^{\mathrm R}$ preserves monomorphisms.
For $i_A:A\hookrightarrow S_1$, a factorization
$k:GA\to B$ with $i_Bk=\sigma G i_A$ transposes to
$k^\flat=G^{\mathrm R}k\,\eta_A:A\to G^{\mathrm R}B$ satisfying
\[
 G^{\mathrm R} i_B\,k^\flat=\sigma^\flat i_A .
\]
Conversely this equation transposes back to the factorization through
$i_B$. By the pullback defining $R(B)$, it follows that
$L(A)\leq B$ exactly when $A\leq R(B)$.

To construct the derivation comparison, suppose
$\sigma_{\Phi}:G\Phi_1\to\Phi_2$ satisfies
$\sigma_\Phi G j_1=j_2\sigma$, where $j_i:S_i\hookrightarrow\Phi_i$.
Define $K_\Phi:\mathcal D_1/\Phi_1\to\mathcal D_2/\Phi_2$ by
$K_{\Phi}(q)=\sigma_{\Phi}Gq$ on objects and $K_\Phi(k)=Gk$ on
slice morphisms. Its right adjoint $R_\Phi$ sends $b:B\to\Phi_2$
to the pullback of $G^{\mathrm R}b$ along
$G^{\mathrm R}\sigma_{\Phi}\,\eta_{\Phi_1}:\Phi_1\to G^{\mathrm R}\Phi_2$.
On a morphism $k:(B,b)\to(B',b')$ over $\Phi_2$, $R_\Phi(k)$
is induced by $1_{\Phi_1}$ and $G^{\mathrm R}k:G^{\mathrm R}B\to G^{\mathrm R}B'$ on these pullbacks.
For each $A$, suppose the comparison of rule instances supplies a
natural transformation
\[
 \lambda:K_{\Phi}H_A\to H_{L(A)}K_{\Phi}.
\]
Its component at $q$ sends a rule instance with premise derivations
over $q$ to a target rule instance with premises over $K_\Phi(q)$
and the translated conclusion. Let
$\delta_{L(A)}:H_{L(A)}P_{L(A)}\to P_{L(A)}$ be the target initial algebra and
$\epsilon:K_{\Phi}R_{\Phi}\to1$ the counit. Transpose
\[
 K_{\Phi}H_AR_{\Phi}P_{L(A)}
 \xrightarrow{\lambda}H_{L(A)}K_{\Phi}R_{\Phi}P_{L(A)}
 \xrightarrow{H_{L(A)}\epsilon}H_{L(A)}P_{L(A)}
 \xrightarrow{\delta_{L(A)}}P_{L(A)}
\]
to give $R_{\Phi}P_{L(A)}$ a source algebra structure. Source initiality
gives a morphism from $P_A$ to this algebra. Its transpose has
underlying morphism $h:GP_A\to P_{L(A)}$ with
$c_{L(A)}h=\sigma_\Phi Gc_A$. Write $k_A:P_{A,S_1}\to P_A$ for the
pullback projection. Then
\[
 c_{L(A)}h\,Gk_A=\sigma_\Phi G j_1\,Gc_{A,S_1}
                 =j_2\sigma Gc_{A,S_1}.
\]
The pullback $P_{L(A),S_2}=P_{L(A)}\times_{\Phi_2} S_2$ therefore gives
$p_A$ with components $h\,Gk_A$ and $\sigma Gc_{A,S_1}$.

\subsection{Free first-order theories and inseparability}
\label{sec:free-first-order-theories}

Assume additionally that every monomorphism of $\mathcal D$ has a
partial inverse in $\mathscr S$. Subobjects then identify with
restriction idempotents. A consistent sentence theory is essentially incomplete if a total morphism
reduces an inseparable pair to its theorem and refutation subobjects
(Corollary~\ref{cor:domain-incompleteness}).

Fix the decidable variable object $V$ of
Appendix~\ref{app:categorical-hypotheses}. For each $X\in\mathcal D$,
take the presentation $P_X$ with $\Pi_0=X$ and all other signature
objects $\Sigma_k,\Pi_k$ initial. Assume that its syntax and
variable-list algebras of Definition~\ref{def:categorical-syntax}
are initial in every parameter slice and that it has an initial
derivation algebra for each axiom subobject of its sentence object
(Appendix~\ref{app:first-order-derivations}).

The terms of $P_X$ are variables. Since $\Tm_X^0=1$, the relation
constructor at arity $0$ in Appendix~\ref{app:first-order-syntax}
gives $r_0:X\to\Phi_X$. Its values are the propositional atoms of $P_X$.
Their free-variable lists are empty, so $r_0$ factors through the
sentence inclusion $j_X:S_X\hookrightarrow\Phi_X$ as
$j_X\theta_X$, with $\theta_X:X\to S_X$.
A generalized element $x:W\to X$ gives the sentence
$\theta_Xx:W\to S_X$. For $h:W'\to W$, the sentence at $xh$ is
$(\theta_Xx)h$.
For disjoint subobjects $a,b\hookrightarrow X$, with inclusions
$i_a,i_b$, put
\[
 A_{a,b}=\im(\theta_Xi_a)\vee\im(\nu_X\theta_Xi_b),
 \qquad T_{a,b}=\operatorname{Gen}_{P_X}(A_{a,b}).
\]

Let $\mathsf{Pair}_c(\mathcal D)$ have triples $(X,a,b)$ as
objects, where $a\wedge b=0$ and $T_{a,b}$ is consistent. A
morphism $f:(X,a,b)\to(Y,c,d)$ is a total $f:X\to Y$ satisfying
$a\leq f^*c$ and $b\leq f^*d$.
Let $\mathsf{FO}_V(\mathcal D)$ have pairs $(P,T)$ as objects,
where $P$ has variable object $V$ and $T$ is a consistent
sentence theory. A morphism $(P,T)\to(Q,U)$ is a logical
translation $\sigma:P\to Q$ with $\sigma_!T\leq U$.
Define
\[
 \Sigma_{\mathrm{FO}}(X,a,b)=(P_X,T_{a,b}),\qquad
 J_{\mathrm{FO}}(P,T)=(S_P,T,\nu_P^*T).
\]
The functor $\Sigma_{\mathrm{FO}}$ sends $f$ to the interpretation
extending $\theta_Yf$, and $J_{\mathrm{FO}}$ sends $\sigma$ to
$\sigma_S$.

Voutsadakis gives an adjunction for relational signatures in which
morphisms substitute formulas for relation symbols
\cite[Theorem~7]{Voutsadakis}. Here an assignment $f:X\to S_P$ of
sentences to propositional atoms extends to a theory morphism
$(P_X,T_{a,b})\to(P,T)$ precisely when $a\leq f^*T$ and
$b\leq f^*\nu_P^*T$.

\begin{theorem}[Free first-order adjunction]
\label{thm:free-first-order-adjunction}
Let $\mathcal D=\Tot(\mathscr S)$ be regular with disjoint,
pullback-stable finite coproducts, where $\mathscr S$ is a split
restriction category and every monomorphism of $\mathcal D$ has a partial
inverse in $\mathscr S$. Fix a decidable variable object $V$. For every
$X\in\mathcal D$, assume that $P_X$ has the syntax and variable-list
algebras of Definition~\ref{def:categorical-syntax}, initial in every
parameter slice, and an initial derivation algebra for each axiom
subobject. In $\mathsf{FO}_V(\mathcal D)$, each presentation $P$
has an initial derivation algebra for every $A\hookrightarrow S_P$,
and $T$ in an object $(P,T)$ is a consistent sentence theory of $P$.
Then
\[
 \Sigma_{\mathrm{FO}}:
 \mathsf{Pair}_c(\mathcal D)\rightleftarrows
 \mathsf{FO}_V(\mathcal D):J_{\mathrm{FO}},\qquad
 \Sigma_{\mathrm{FO}}\dashv J_{\mathrm{FO}}.
\]
The unit at $(X,a,b)$ has underlying morphism $\theta_X:X\to S_X$.
The counit $\Sigma_{\mathrm{FO}}J_{\mathrm{FO}}(P,T)\to(P,T)$
is the translation extending $1_{S_P}:S_P\to S_P$.
\end{theorem}
\begin{proof}
Let $f:X\to S_P$ assign a sentence to each propositional atom, and
interpret variables by $V\to\Tm_P$. Initiality of the free syntax
gives a unique term and formula interpretation $\widehat f$.
Let $\ell:W_X\xrightarrow{\sim}W_P$ be the unique isomorphism
of the two initial variable-list algebras. Write
$\FV_X:\Phi_X\to W_X$ and $\FV_P:\Phi_P\to W_P$ for the
free-variable occurrence morphisms. Since $f$ takes values in
sentences, the interpretation of each propositional atom has an empty
occurrence list.
At a quantifier, $\ell$ commutes with deletion of the bound variable;
at equality and implication it commutes with concatenation.
The occurrence folds therefore give
$\FV_P\widehat f_\Phi=\ell\FV_X$ on formulas and
$\FV_{\Tm_P}\widehat f_{\Tm}=\ell\FV_{\Tm_X}$ on terms.
Pullback along the empty-list inclusion gives the sentence morphism $\widehat f_S:S_X\to S_P$.

For substitution of $t$ for $x$ in $\forall y\,\varphi$, the branch
$y=x$ retains the quantified formula. On the branch $y\ne x$,
membership in the occurrence lists is preserved by $\ell$.
Hence the two interpretations give the same values for the tests
$x\in\FV_X(\varphi)$ and $y\in\FV_{\Tm}(t)$ that determine
capture in~\eqref{eq:quantifier-fold}. At a propositional atom,
substitution fixes its interpreting sentence by the absent-variable
identity of Appendix~\ref{app:substitution}. At an equality $u=v$,
substitution acts on the two terms. Put $E_X=V\times\Tm_X$,
$E_P=V\times\Tm_P$ and $e=1_V\times\widehat f_{\Tm}:E_X\to E_P$.
Write $b_X:\Phi_X\times E_X\to 1+\Phi_X$ and
$b_P:\Phi_P\times E_P\to 1+\Phi_P$ for formula substitution with
failure in the first summand. The constructor equations and initiality
over $E_X$ give
\[
 (1_1+\widehat f_\Phi)b_X=b_P(\widehat f_\Phi\times e).
\]
Pulling back the success summand gives $s:\Adm_X\to\Adm_P$ with
the two substitution equations above. These equations translate the
instantiation and equality-substitution instances, and preservation
of the free-variable test translates quantifier distribution.
The other four logical schemes are preserved by their constructors.
Thus $\widehat f:P_X\to P$ is a logical translation.

Take $X=S_P$ and $f=1_{S_P}$, and denote the resulting translation
by $\epsilon:P_{S_P}\to P$. The theory $\epsilon_S^*T$ is closed by
Theorem~\ref{thm:translation-galois}. It is consistent because
$\epsilon_S$ commutes with negation and its inverse-image map preserves
meets and bottom. As $\epsilon_S\theta_{S_P}=1_{S_P}$, it contains
the positive axioms specified by $T$ and the negative axioms specified
by $\nu_P^*T$. Their generated theory is therefore consistent.
This proves that $J_{\mathrm{FO}}(P,T)$ is an object of
$\mathsf{Pair}_c(\mathcal D)$. For a theory morphism
$\sigma:(P,T)\to(Q,U)$, the bound $T\leq\sigma_S^*U$ and
$\sigma_S\nu_P=\nu_Q\sigma_S$ give
$\nu_P^*T\leq\sigma_S^*\nu_Q^*U$. Thus $\sigma_S$ is a morphism of pairs.

For a pair morphism $f:(X,a,b)\to J_{\mathrm{FO}}(P,T)$, the
defining inequalities put $\widehat f_!A_{a,b}$ inside $T$.
Theorem~\ref{thm:translation-galois} gives
$\widehat f_!T_{a,b}\leq T$.

Conversely, a theory morphism $\sigma:(P_X,T_{a,b})\to(P,T)$
restricts to $f=\sigma_S\theta_X:X\to S_P$. Preservation of negation
and the definition of the signed axioms give
\[
 a\leq f^*T,\qquad b\leq f^*\nu_P^*T.
\]
The extension satisfies $\widehat f_S\theta_X=f$, and uniqueness gives
$\widehat{\sigma_S\theta_X}=\sigma$. These equations make the two
assignments inverse. For a morphism of pairs $k$ with underlying
morphism $X'\to X$ and a theory translation $\rho:(P,T)\to(Q,U)$,
both sides of
\[
 \widehat{\rho_S f k}=\rho\,\widehat f\,\Sigma_{\mathrm{FO}}(k)
\]
take each propositional atom to the same sentence in $Q$.
Uniqueness of extension gives the equation, proving naturality in both
variables. The same uniqueness gives the functor laws for
$\Sigma_{\mathrm{FO}}$. For $J_{\mathrm{FO}}$ they follow from
composition of sentence morphisms.
\end{proof}

\begin{corollary}\label{cor:signed-first-order-incompleteness}
Under the hypotheses of
Theorem~\ref{thm:free-first-order-adjunction}, let $a,b$ be
disjoint subobjects of $X$ and suppose $T_{a,b}$ is consistent.
Then $T_{a,b}$ is essentially incomplete if and only if $(a,b)$
is inseparable in $\mathscr S$, that is, if and only if no complementary
subobjects $C,D$ of $X$ satisfy $a\leq C$ and $b\leq D$.
\end{corollary}
\begin{proof}
Let $U\geq T_{a,b}$ be any consistent closed theory of $P_X$.
Axiom introduction and monotonicity of inverse image give
\[
 a\leq\theta_X^*T_{a,b}\leq\theta_X^*U,\qquad
 b\leq\theta_X^*\nu_X^*T_{a,b}\leq\theta_X^*\nu_X^*U .
\]
If $U$ were complete, $U$ and $\nu_X^*U$ would be complementary.
Their inverse images under $\theta_X$ would separate $a,b$, as in
Corollary~\ref{cor:domain-incompleteness}. Thus inseparability
implies essential incompleteness.

Conversely, let $m:C\hookrightarrow X$ and $n:D\hookrightarrow X$
be complementary subobjects containing $a$ and $b$, respectively.
Disjointness and extensivity make $[m,n]:C+D\to X$ monic, and
its image is the join of the two subobjects, hence all of $X$.
Thus $[m,n]$ is an isomorphism. Let $\Two=1+1$ with injections
$\mathbf{f},\mathbf{t}:1\to\Two$, in that order. The total morphism
$d=[\mathbf{t}!_C,\mathbf{f}!_D][m,n]^{-1}:X\to\Two$
is true on $a$ and false on $b$.
Interpret the terms in $1$, the propositional atoms by $d$, equality
by $\mathbf{t}:1\to\Two$ under $1\times1\cong1$, negation and
implication by their Boolean truth functions, and universal quantification
by the projection
$\pi_2:V\times\Two\to\Two$.
Free syntax gives a total valuation $v:\Phi_X\to\Two$.
In the substitution fold of Appendix~\ref{app:substitution}, consider
the pairs of an original formula and a substitution result. In
$\Phi_X\times(1+\Phi_X)$ take the coproduct of the failure summand
$\Phi_X\times1$ and the equalizer of $v\pi_1,v\pi_2:\Phi_X^2\to\Two$
in the success summand.
The atomic results lie in this subobject because substitution fixes every
propositional atom and every equality has value $\mathbf t$. Negation and implication
preserve equality of values whenever their substituted arguments succeed.
At $\forall y\,\varphi$, substitution for $x$ returns the original
formula when $y=x$. When $y\ne x$, a successful result has the value
of the substituted body, hence of $\varphi$. Every failure belongs to
the first summand. Thus the subobject, with the full term carrier, is
closed under the substitution algebra in the parameter slice.
Initiality puts the fold in this subobject. On the success summand
$i:\Adm\hookrightarrow\Phi_X\times V\times\Tm_X$, with result
$a_F:\Adm\to\Phi_X$, this gives $v a_F=v\pi_{\Phi_X}i$.
Thus instantiation and
substitutivity of equality have value
$\mathbf{t}$. Reflexivity has that value by the interpretation of
equality, and quantifier distribution has it because quantification leaves
each truth value unchanged. The propositional schemes have value
$\mathbf{t}$, and modus ponens and generalization preserve truth.
Thus $U=(v j_X)^*\mathbf{t}\hookrightarrow S_X$ is closed. Negation
interchanges true and false, so $\nu_X^*U=(v j_X)^*\mathbf{f}$.
The two fibers are complementary, proving consistency and completeness.
Since $v j_X\theta_X=d$, the signed axioms lie in $U$, whence
$T_{a,b}\leq U$.
This contradicts essential incompleteness.
\end{proof}
